\documentclass[11pt,a4paper]{article}
\usepackage[utf8]{inputenc}
\usepackage[T1]{fontenc}
\usepackage[english]{babel}
\usepackage{lmodern}
\usepackage{geometry}
\usepackage{xcolor}
\usepackage{tocloft}
\usepackage{hyperref}
\usepackage{enumitem}
\usepackage{parskip}
\usepackage{booktabs}
\usepackage{longtable}
\usepackage{fancyhdr}
\usepackage{lastpage}
\usepackage{listings}
\usepackage{titlesec}
\geometry{margin=2.15cm}
\setlength{\headheight}{14pt}
\setlength{\emergencystretch}{3em}
\setlist{itemsep=3pt,parsep=1pt,topsep=4pt}
\definecolor{titleblue}{HTML}{123B5D}
\definecolor{softblue}{HTML}{EDF5FB}
\definecolor{softgray}{HTML}{F5F7FA}
\definecolor{accent}{HTML}{0F6E8C}
\hypersetup{
  colorlinks=true,linkcolor=titleblue,urlcolor=accent,
  pdftitle={Scientific Data Analysis v2.6: English publication edition},
  pdfauthor={interemi},
  pdfsubject={Measured usage and conservative cost optimization}
}
\titleformat{\section}{\Large\bfseries\color{titleblue}}{\thesection.}{0.6em}{}
\titleformat{\subsection}{\large\bfseries\color{accent}}{\thesubsection.}{0.6em}{}
\setlength{\cftsecnumwidth}{2.8em}
\tocloftpagestyle{fancy}
\lstdefinestyle{promptstyle}{
  basicstyle=\ttfamily\small,backgroundcolor=\color{softgray},
  frame=single,rulecolor=\color{softblue},breaklines=true,columns=fullflexible
}
\pagestyle{fancy}
\fancyhf{}
\fancyhead[L]{Scientific Data Analysis}
\fancyhead[R]{v2.6 guide: English edition}
\fancyfoot[C]{\thepage\ of \pageref{LastPage}}
\newcommand{\code}[1]{\texttt{\detokenize{#1}}}

\begin{document}
\pagenumbering{gobble}
\begin{titlepage}
\centering
\vspace*{1.0cm}
{\Huge\bfseries\color{titleblue} Measured Usage Guide\\[0.15cm]
Scientific Data Analysis v2.6\par}
\vspace{0.65cm}
{\Large Reproducible observation and evidence-based optimization\par}
\vspace{0.8cm}
{\large Operational health: benchmark matrix, local harness,
mother/child comparison, and app-readiness gates\par}
\vspace{0.8cm}
{\large English publication edition: September 22, 2026\par}
\vspace{0.5cm}
\begin{minipage}{0.93\textwidth}
\small\raggedright
Translation of the preserved private v2.6 guide at commit \code{d6583f1}.
Original TeX SHA-256:
\nolinkurl{86bf78142300618242b3036d4537e88009ac353d30b96a798d4c2d66a562c131}.
Private source:
\nolinkurl{skills/scientific-data-analysis/.cache/active_references/guides/guia_scientific_data_analysis_v2_6.tex}.

Recorded results and closure steps are historical, not current verification.
This release guide does not synchronize or modify Scientific Workbench;
the app remains outside its scope. Historical commands are retained for
traceability: use fresh output directories when reproducing them, without
overwriting earlier evidence. The original remains unchanged privately.
\end{minipage}
\vfill
{\large Prepared with Codex\par}
{\large Copyright 2026 interemi\par}
\end{titlepage}

\clearpage
\pagenumbering{roman}
\fancyfoot[C]{\thepage}
\begin{abstract}
\code{scientific-data-analysis} v2.6 continues the modular architecture
stabilized in v2.5, with a different question. v2.5 reduced static loading and
closed required plugin-eval fixes. v2.6 asks where representative usage
actually incurs cost. It introduces a benchmark matrix, an observed-usage
harness, reproducible \code{observed_usage.jsonl} results, a mother/child
comparison, and a conservative policy: optimize only when measurements
justify the change.

The measurement is a local proxy, not a reading of actual weekly Codex tokens.
v2.6 therefore promises no reduction in account usage limits. It provides
verifiable infrastructure for measuring and comparing future decisions while
preserving v2.5 compatibility.
\end{abstract}
\tableofcontents
\newpage
\pagenumbering{arabic}
\fancyfoot[C]{\thepage\ of \pageref{LastPage}}

\section{Executive overview}
v2.6 focuses on measurement and restraint. It adds no capabilities, changes no
scientific contracts, and leaves Scientific Workbench untouched. It makes
efficiency discussions evidence-based through reproducible scenarios,
parseable outputs, and an initial ranking of local cost.

The product decisions are:
\begin{itemize}[leftmargin=1.5em]
\item \code{scientific-data-analysis} remains the normal entrypoint for ambiguous or mixed tasks.
\item Use \code{scientific-data-astro}, \code{scientific-data-documents}, \code{scientific-data-notebooks}, or \code{scientific-data-maintainer} when the intent is already specific.
\item Results use \code{token_measurement_kind=local_proxy}; they are not actual account token counts.
\item Registries, artifact types, run bundles, app hints, and app-readiness contracts remain unchanged.
\item Optimization is limited to stabilizing the harness and documenting measurable decisions.
\end{itemize}

\section{From v2.5 to v2.6}
\begin{longtable}{p{0.20\linewidth}p{0.68\linewidth}}
\toprule
\textbf{Area} & \textbf{v2.6 change} \\
\midrule
Charter & \code{v2-6-measured-usage-charter.md} defines objectives, non-goals, and closure criteria. \\
Observed contract & \code{v2-6-observed-usage-contract.md} separates local proxy measurement from actual tokens. \\
Benchmark matrix & \code{v2-6-benchmark-matrix.json} contains 8 synthetic or anonymized scenarios. \\
Harness & \code{audit_v2_6_observed_usage_harness.py} runs scenarios and saves \code{observed_usage.jsonl}. \\
Regressions & Adds v2.6 checks for baseline, matrix, observed usage, and app-readiness sanity. \\
Router comparison & \code{v2-6-routing-cost-findings.md} documents mother versus child routes. \\
Optimization & \code{v2-6-cost-optimization-log.md} records decisions; no large refactor or public trigger change. \\
Maintainer scenario & Uses a maintainer-specific regression and captures JSON stdout as a parseable summary. \\
Guide & Replaces v2.5 as the active release document. \\
\bottomrule
\end{longtable}

\section{What remains unchanged}
\begin{itemize}[leftmargin=1.5em]
\item No new capability, public path deletion, or public path rename.
\item No edits to Scientific Workbench.
\item No synchronization of the installed family until the final phase.
\item No forced restructuring of mother and children.
\item No claim of actual token savings without comparable actual measurements.
\item No changes to the scientific contracts of astro, documents, notebooks, or maintainer routes.
\end{itemize}

\section{Measured scenarios}
The v2.6 matrix covers eight scenarios. All inputs are synthetic or anonymized;
originals are not modified.

\begin{longtable}{p{0.43\linewidth}p{0.45\linewidth}}
\toprule
\textbf{Scenario} & \textbf{Owner, route, and input} \\
\midrule
\code{table_professional_csv} & notebooks; direct child; synthetic professional CSV \\
\code{mixed_folder_router} & mother; router; synthetic mixed folder \\
\code{legacy_notebook_inspect} & notebooks; direct child; synthetic legacy notebook \\
\code{document_admin_intake} & documents; direct child; administrative document package \\
\code{fits_simple_inspect} & astro; direct child; minimal synthetic FITS \\
\code{optional_stilts_block} & astro; optional; simulated absent STILTS backend \\
\code{maintainer_surface_check} & maintainer; maintainer route; v2.5 maintainer regression \\
\code{router_vs_direct_table} & mother; router; the same CSV used for direct profiling \\
\bottomrule
\end{longtable}

\section{Observed result}
Reference execution:
\begin{lstlisting}[style=promptstyle]
python scripts/audit_v2_6_observed_usage_harness.py \
  --output-dir tmp/v2_6_phase2_observed_usage_harness
\end{lstlisting}

\begin{longtable}{p{0.41\linewidth}p{0.23\linewidth}rr}
\toprule
\textbf{Scenario} & \textbf{Status} & \textbf{Proxy} & \textbf{ms} \\
\midrule
\code{table_professional_csv} & PASS & 1975 & 459 \\
\code{mixed_folder_router} & PASS & 20484 & 94 \\
\code{legacy_notebook_inspect} & Controlled block & 811 & 67 \\
\code{document_admin_intake} & PASS & 3605 & 602 \\
\code{fits_simple_inspect} & PASS & 1313 & 110 \\
\code{optional_stilts_block} & Controlled block & 1457 & 73 \\
\code{maintainer_surface_check} & PASS & 162 & 44 \\
\code{router_vs_direct_table} & PASS & 20565 & 90 \\
\bottomrule
\end{longtable}
``Proxy'' is the local token proxy; ``ms'' is duration in milliseconds.
``Controlled block'' preserves the original machine status
\code{BLOCKED_CONTROLADO}.

Interpretation:
\begin{itemize}[leftmargin=1.5em]
\item There are 6 PASS results and 2 controlled blocks.
\item The notebook stops cleanly under system Python because \code{nbformat} is missing.
\item The absent STILTS backend blocks cleanly; this is not treated as a core error.
\item The CSV comparison shows a heavier mother route than the direct child when the goal is already specific.
\item \code{local_proxy} supports internal comparisons, not promises of actual weekly savings.
\end{itemize}

\section{Mother versus child routing}
The mother remains the appropriate entrypoint for broad requests such as
``inspect this folder'' or ``help me with this project.'' Direct children
provide a shorter route when the intent is clear:
\begin{itemize}[leftmargin=1.5em]
\item FITS, WCS, photometry, spectra, or astronomy: \code{scientific-data-astro}.
\item PDF, DOCX, LaTeX, presentations, or reporting: \code{scientific-data-documents}.
\item CSV, notebooks, containers, time series, or tables: \code{scientific-data-notebooks}.
\item Release gates, plugin-eval, or maintenance: \code{scientific-data-maintainer}.
\end{itemize}
v2.6 does not change trigger text or public contracts on the basis of one
measurement. It documents evidence and provides repeatable regression checks.

\section{v2.6 regressions}
Main regressions:
\begin{lstlisting}[style=promptstyle]
python scripts/audit_v2_6_phase0_baseline_regression.py
python scripts/audit_v2_6_benchmark_matrix_regression.py
python scripts/audit_v2_6_observed_usage_harness.py --output-dir tmp/v2_6_phase2_observed_usage_harness
python scripts/audit_v2_6_observed_usage_regression.py --output-dir tmp/v2_6_phase2_observed_usage_harness
python scripts/audit_v2_6_app_ready_cost_sanity_regression.py --output-dir tmp/v2_6_phase5_app_ready_sanity
\end{lstlisting}
Inherited gates that must continue to pass:
\begin{lstlisting}[style=promptstyle]
python scripts/audit_v2_5_budget_deferral_regression.py
python scripts/audit_v2_4_budget_architecture_regression.py
python scripts/sync_public_surface_docs.py --check
python scripts/portable_smoke_test.py --profile core
\end{lstlisting}

\section{Historical closure criteria}
The original v2.6 closure required all ten conditions:
\begin{enumerate}[leftmargin=1.7em]
\item The benchmark matrix covers all 5 skills.
\item The harness produces \code{observed_usage.jsonl}, \code{result.json}, and \code{summary.md}.
\item All 8 scenarios are PASS, WARNING, or justified \code{BLOCKED_CONTROLADO}, with no FAIL/ROTO.
\item App-readiness sanity validates summaries, artifacts, statuses, and \code{original_modified=false}.
\item \code{guia_scientific_data_analysis_v2_6.tex} compiles to PDF.
\item The PDF and TeX are copied to the Desktop, as required by that historical local handoff.
\item v2.6 and inherited gates pass in the editable family.
\item Installed synchronization occurs only after green gates.
\item Post-synchronization gates pass.
\item \code{evaluate-skill} runs against the installed mother and children.
\end{enumerate}

\section{Recommended usage after v2.6}
For broad tasks:
\begin{lstlisting}[style=promptstyle]
python scripts/scientific_workflow_router.py plan input --task "..."
\end{lstlisting}
For specific tasks:
\begin{lstlisting}[style=promptstyle]
python scripts/profile_table.py data.csv --summary-json summary.json
python scripts/inspect_fits.py image.fits --summary-json summary.json
python scripts/document_intake_workbench.py docs/ --summary-json summary.json
\end{lstlisting}
To repeat the v2.6 measurement:
\begin{lstlisting}[style=promptstyle]
python scripts/audit_v2_6_observed_usage_harness.py --output-dir tmp/v2_6_observed_usage
python scripts/audit_v2_6_observed_usage_regression.py --output-dir tmp/v2_6_observed_usage
\end{lstlisting}

\section{Plugin-eval evaluation update}
After conservative hardening following v2.6, the mother and four children were
optimized for \code{plugin-eval:evaluate-skill} without removing capabilities
or breaking public paths. The full content was deferred: guides, long
references, historical duplicates, examples, and heavy helpers remained
available outside the initial content Codex reads when invoking a skill.

The remaining informational item was \code{coverage-artifacts-unavailable}.
To address it, \code{audit_v2_7_plugin_eval_coverage_artifacts.py} generates
\code{coverage-summary.json} for each skill. It reports smoke/regression
coverage observed with \code{trace}, not exhaustive project coverage. It gives
\code{plugin-eval} a real, traceable coverage signal instead of an empty field.

Recorded editable-family state before synchronization:
\begin{longtable}{p{0.43\linewidth}p{0.18\linewidth}p{0.25\linewidth}}
\toprule
\textbf{Skill} & \textbf{Rating} & \textbf{Deductions} \\
\midrule
\code{scientific-data-analysis} & 100/A & none \\
\code{scientific-data-astro} & 100/A & none \\
\code{scientific-data-documents} & 100/A & none \\
\code{scientific-data-notebooks} & 100/A & none \\
\code{scientific-data-maintainer} & 100/A & none \\
\bottomrule
\end{longtable}
Three invariants must remain:
\begin{itemize}[leftmargin=1.5em]
\item Public paths continue to exist, even when they point to deferred content.
\item Coverage artifacts are small, parseable, and regenerable.
\item \code{evaluate-skill} runs on mother and children after every synchronization.
\end{itemize}

\section{Conclusion}
v2.6 provides reproducible evidence for reasoning about cost, preserves the
v2.5 architecture, and avoids unnecessary contract changes. Use the mother for
ambiguity, children for defined tasks, and require benchmark evidence before
future optimizations.
\end{document}
